\section{Conclusion}
\label{sec:conclusion}

Cette première phase d'expérimentation a permis d'établir avec succès une chaîne complète de tomodensitométrie à rayons X pour l'inspection volumique non destructive de pièces issues de la fabrication additive. En mobilisant un banc d'imagerie de laboratoire et un protocole de balayage haute résolution à $720$ projections sur $360^\circ$, les objectifs scientifiques initiaux ont été pleinement atteints.

L'analyse systématique des architectures cinématiques a permis de résoudre l'ambiguïté d'alignement de l'axe de rotation mécanique et de démontrer que la découpe le long des colonnes du capteur (Famille 2) est l'unique configuration fidèle à la géométrie physique du système. La calibration sub-pixel du centre de rotation a révélé un décalage optimal de $\Delta x_{\text{COR}} = -11{,}6 \pm 0{,}2~\text{pixels}$, dont la compensation a éliminé les artéfacts de dédoublement et conféré un gain de rapport signal sur bruit de $+14{,}2~\text{dB}$.

Sur le plan métrologique, l'extraction de profils d'atténuation et la mesure à mi-hauteur (FWHM) ont permis de quantifier l'épaisseur de la coquille extérieure du modèle en PLA à $0{,}81 \pm 0{,}04~\text{mm}$, en accord parfait ($+1{,}25\,\%$, inférieur à l'incertitude à $1\,\sigma$) avec la valeur nominale de consigne ($0{,}80 \pm 0{,}02~\text{mm}$). De surcroît, le réseau interne d'alvéoles de remplissage a été visualisé et mesuré ($3{,}18 \pm 0{,}08~\text{mm}$), confirmant la sensibilité de la tomodensitométrie aux phénomènes de retrait thermique post-extrusion.

Ces résultats valident sans ambiguïté les capacités du dispositif et préparent directement la seconde étape de l'étude. Dès la prochaine séance, ce protocole optimisé sera appliqué aux scans tomographiques de fantômes cubiques étalonnés intégrant des canaux et des cavités de dimensions micrométriques certifiées. Cette suite logique permettra de clore la boucle métrologique en dressant une cartographie quantitative des écarts dimensionnels réels par rapport aux fichiers CAO d'origine.
